\section{Numerical error summaries}
\quad This appendix collects complementary NRMSE summaries referenced in Sec. V. These plots supplement the representative scatter comparisons shown in the main text by providing ensemble-level error statistics across multiple Erd\H{o}s-R\'enyi network realizations, as well as compact summaries of how the first-order prediction becomes less accurate as the number or magnitude of simultaneous perturbations increases.

\begin{figure}[H]
    \centering
    \includegraphics*[width=.7\columnwidth]{Figure/Else/NRMSE_Bar_Chart.eps}
    \caption{Average NRMSE for the single-parameter steady-state response tests under three background conditions: equilibrium (EQ), nonequilibrium steady states generated by directed couplings (NESS-Dir), and nonequilibrium steady states generated by heterogeneous noise strengths (NESS-Noise). For each background, the node category corresponds to perturbations of a single node parameter and the edge category corresponds to perturbations of a single connection weight. Each marker shows the mean NRMSE obtained from 20 independent Erd\H{o}s-R\'enyi network configurations, and the error bars indicate one standard deviation across these realizations. The small averaged errors across all six cases support the robustness of the first-order response prediction beyond the representative examples shown in the main text.}
    \label{app:NRMSE_Bar_Chart}
\end{figure}

\begin{figure}[H]
    \centering
    \subfigure[]{\includegraphics*[width=.45\columnwidth]{Figure/Else/NRMSE_Node_Perturb.eps}}
    \subfigure[]{\includegraphics*[width=.46\columnwidth]{Figure/Else/NRMSE_Edge_Perturb.eps}}
    \caption{Average NRMSE of the superposed linear-response prediction for multiple simultaneous perturbations in an undirected Erd\H{o}s-R\'enyi network (\(N=100\), \(p=0.1\)) with nonuniform noise strengths. Panel (a) shows the mean NRMSE as a function of the number of perturbed nodes, and panel (b) shows the corresponding mean NRMSE as a function of the number of perturbed undirected edges. For each data point, the NRMSE is averaged over 20 independent network configurations; within each realization, the perturbed nodes or edges are selected at random. The perturbation amplitudes are fixed at \(\delta r_{\alpha}=0.5\) and \(\delta W_{\alpha\beta}=0.5\), which are approximately one tenth of the mean sampled values of \(r\) and \(W\). The small averaged NRMSE over a wide range indicates that the first-order response remains quantitatively useful for substantial but still moderate multiple perturbations.}
    \label{app:NRMSE_Error_Perturb}
\end{figure}

\begin{figure}[H]
    \centering
    \subfigure[]{\includegraphics*[width=.45\columnwidth]{Figure/Else/NRMSE_Edge_Remove.eps}}
    \subfigure[]{\includegraphics*[width=.45\columnwidth]{Figure/Else/NRMSE_Edge_Weaken.eps}}
    \caption{Mean NRMSE of the first-order response prediction for removal-type perturbations in an undirected Erd\H{o}s-R\'enyi network (\(N=100\), \(p=0.1\)) with nonuniform noise strengths. Panel (a) shows the mean NRMSE as a function of the number of removed undirected edges. For each data point, the NRMSE is averaged over 20 independent network configurations, with the removed edges chosen at random in each realization. Panel (b) shows the corresponding mean NRMSE for node weakening, where a node is chosen at random and all edges incident to it are scaled according to \(\delta W_{\alpha\beta}=-\lambda W_{\alpha\beta}\); \(\lambda=1\) corresponds to complete node removal. The growing error demonstrates the breakdown of the linear approximation under sufficiently large structural perturbations.}
    \label{app:NRMSE_Error_Remove}
\end{figure}
